
\chapter{素数分布：从古典经验公式到解析理论}\label{ch:prime_distribution}

\vspace{0.6em}
\begin{quote}
\small
我很快看出：在种种起伏之下，素数出现的频率平均与对数成反比，
因而小于给定界限 $n$ 的素数个数近似等于
$\int \frac{\mathrm{d}n}{\log n}$。\\[4pt]
\hfill——C.~F.~Gauss，致 J.~F.~Encke 的信（1849 年 12 月 24 日）
\end{quote}
\vspace{0.6em}

素数是大于 $1$、且仅有正因子 $1$ 与自身的正整数；大于 $1$ 的整数可唯一分解为素数之积（算术基本定理）、
素数有无穷多个（欧几里得）。
素数在数轴上的排列方式，以及给定范围内的个数，是由上述乘数因子唯一分解所引出的基本分布问题。
Riemann 将此问题置于复分析框架，为 Hadamard 与 de~la~Vallée-Poussin
于 1896 年独立完成的素数定理证明提供了理论基础。

本章给出分布问题、古典经验公式与素数定理的\textbf{表述}
（证明见后文），并预示后文解析路线的主要线索：
Euler 乘积、非平凡零点与显式公式的形状。
素数定理的证明提要见第~\ref{ch:prime_number_theorem}~章（细节见附录~\ref{app:pnt_estimates}）；
显式公式见第~\ref{ch:riemann_explicit_formula}~章；
Riemann 1859 年原稿解读见第~\ref{ch:riemann_original_interpretation}~章。

%===========================================================
\section{问题的提出：\texorpdfstring{$\pi(x)$}{pi(x)}}\label{sec:prime_atoms}
%===========================================================

\begin{definition}{素数计数函数}{}
\label{def:pi-ch1}
对实数 $x \geq 2$，定义
\[
\pi(x)
=\#\{p \leq x : p \;\text{为素数}\}
=\sum_{\substack{p \leq x \\ p\,\text{素数}}} 1,
\]
即不超过 $x$ 的素数个数。
两种写法等价：左端是集合基数，右端是对每个素数贡献 $1$ 的求和；
第~\ref{ch:elementary_number_theoretic_functions}~章沿用和式，并系统给出 $\pi,J,\psi$ 等阶梯函数。
\end{definition}

$\pi(x)$ 是右连续的阶梯函数：在每个素数 $p$ 处增加 $1$，在相邻素数之间取常值。
分布问题的定量形式是：当 $x\to\infty$ 时，$\pi(x)$ 大约如何增长？与光滑逼近之差如何控制？

%===========================================================
\section{古典经验公式：Legendre 型近似与 Gauss \texorpdfstring{$\Li$}{Li}}\label{sec:classical_main_terms}
%===========================================================

在素数定理证明之前，对 $\pi(x)$ 的光滑逼近属于\textbf{经验公式}：
依据素数表的数值对照提出，用以刻画不超过 $x$ 的素数个数。

\begin{itemize}
  \item \textbf{Legendre 型近似}：
    A.~M.~Legendre 在 \textit{Essai sur la th\'{e}orie des nombres}
    （约 1798 年，1808 年修订）中，根据当时的素数计数，提出
    \[
      \pi(x)\;\approx\;\frac{x}{\log x-A},
    \]
    并取拟合常数 $A\approx 1.08366$（历史上常写 $\ln$；本书 $\log$ 为自然对数）。
    最简情形 $A=0$ 即 $x/\log x$，用以描述 $\pi(x)$ 的大致增长；
    $A\neq 0$ 时，$A$ 由数值拟合得到，用以改进 $x/\log x$ 的逼近。
    其与 $\Li$ 展开中高阶项的关系，见下文~\eqref{eq:Li-asymptotic} 及分母展开~\eqref{eq:legendre-denom-expand}。
  \item \textbf{Gauss 对数积分}：
    C.~F.~Gauss 早年根据素数表观察，提出不超过 $n$ 的素数个数近似等于
    $\int \mathrm{d}n/\log n$
    （1849 年 12 月 24 日致 J.~F.~Encke 的信中重述）。
    本书采用偏移对数积分
    \begin{equation}
    \Li(x) = \int_2^x \frac{\mathrm{d}t}{\log t}, \qquad x > 1
    \label{eq:Li-def-parts}
    \end{equation}
    作为实函数 $\Li:(1,+\infty)\to\mathbb{R}$，是本章与全书默认的光滑主项写法
    （定义域为实上限；复对数积分见第~\ref{ch:logint_Ei}~章）。
    不同于 Legendre 在分母 $\log x-A$ 中插入拟合常数 $A$，
    $\Li$ 以积分直接给出光滑的经验公式，无需额外参数。
    （$\li$、$\Ei$ 与「PNT 用 $\Li$、显式公式用 $\li$」的完整论述见第~\ref{ch:logint_Ei}~章；
    本章先用实 $\Li$ 完成分部展开与数值对照。）
\end{itemize}

二者均为以连续函数逼近阶梯 $\pi(x)$ 的经验公式，
当时既不给出误差的精确结构，也不能证明 $\pi\sim\Li$ 或 $\pi\sim x/\log x$。
下面用分部积分说明：$x/\log x$ 恰为 $\Li(x)$ 渐近展开的\textbf{首项}，
故 $\Li$ 较单一的 $x/\log x$ 更精细；
与 $\pi$ 的渐近关系则由素数定理给出（表述见下节，证明见后文）。
后文黎曼显式公式（第~\ref{ch:riemann_explicit_formula}~章）中，
$\li$ 或 $\Li$ 才作为与零点振荡相对的\textbf{主项}出现。

\subsection{\texorpdfstring{$\Li(x)$}{Li(x)} 的分部积分与渐近展开}
\label{sec:Li-integration-by-parts}

对 $\int\mathrm{d}t/\log t$ 作分部积分：先推导 $t/(\log t)^{m}$ 的微分恒等式，
再两端积分；此写法与经典分部积分等价。
对 $t>1$，
\[
  \frac{\mathrm{d}}{\mathrm{d}t}\Biggl(\frac{t}{\log t}\Biggr)
  =\frac{1}{\log t}-\frac{1}{(\log t)^{2}},
\]
从而
\[
  \frac{1}{\log t}
  =\frac{\mathrm{d}}{\mathrm{d}t}\Biggl(\frac{t}{\log t}\Biggr)
  +\frac{1}{(\log t)^{2}}.
\]
两端从 $2$ 积分到 $x$，得
\begin{align}
  \Li(x)
  &=\Biggl[\frac{t}{\log t}\Biggr]_{2}^{x}
  +\int_{2}^{x}\frac{\mathrm{d}t}{(\log t)^{2}}
  \nonumber\\
  &=\frac{x}{\log x}-\frac{2}{\log 2}
  +\int_{2}^{x}\frac{\mathrm{d}t}{(\log t)^{2}}.
  \label{eq:Li-parts-1}
\end{align}
故 $x/\log x$ 是 $\Li(x)$ 渐近展开的首项。
右端积分项恒正，常数 $-2/\log 2$ 为负；当 $x$ 充分大时积分占主导，于是
\[
  \Li(x) > \frac{x}{\log x}.
\]
首项之后的余项 $\int_2^x\mathrm{d}t/(\log t)^2$ 与常数修正使 $\Li(x)$ 大于 $x/\log x$。

一般地，对整数 $m\ge 1$，
\[
  \frac{1}{(\log t)^{m}}
  =\frac{\mathrm{d}}{\mathrm{d}t}\Biggl(\frac{t}{(\log t)^{m}}\Biggr)
  +\frac{m}{(\log t)^{m+1}}.
\]
反复代入~\eqref{eq:Li-parts-1}，对任意正整数 $n$ 有
\begin{equation}
  \Li(x)
  =\sum_{k=1}^{n}(k-1)!
  \Biggl(
    \frac{x}{(\log x)^{k}}-\frac{2}{(\log 2)^{k}}
  \Biggr)
  +n!\int_{2}^{x}\frac{\mathrm{d}t}{(\log t)^{n+1}},
  \label{eq:Li-parts-finite}
\end{equation}
其中 $k=1$ 时 $0!=1$。余项
$R_n(x)=n!\int_2^x\mathrm{d}t/(\log t)^{n+1}$
满足 $R_n(x)\sim n!\,x/(\log x)^{n+1}$（$x\to+\infty$），
比第 $n$ 项高一阶小。因此
\begin{equation}
  \boxed{
  \Li(x)
  \sim
  \sum_{k=1}^{\infty}(k-1)!\,\frac{x}{(\log x)^{k}}
  }
  \label{eq:Li-asymptotic}
\end{equation}
即
\[
  \Li(x)
  \sim
  \frac{x}{\log x}
  +\frac{x}{(\log x)^{2}}
  +\frac{2x}{(\log x)^{3}}
  +\frac{6x}{(\log x)^{4}}
  +\cdots.
\]
这是渐近级数（一般不收敛）；实用时取有限截断，误差由下一项量级控制。
当 $x$ 充分大时，首项之后各项均为正，故再次有 $\Li(x)>x/\log x$，
高出部分来自 $x/(\log x)^2+2x/(\log x)^3+\cdots$。
下限常数 $-\sum_{k=1}^{n}(k-1)!\,2/(\log 2)^{k}$ 在写 $x\to\infty$ 的渐近主部时并入 $O(1)$。
后文专章给出 $\li$、$\Ei$ 及 $\Li(x)=\li(x)-\li(2)$（$x>1$ 实）
（第~\ref{ch:logint_Ei}~章）；本章展开对实 $\Li$ 进行，与 $\li$ 渐近主部相同。

\paragraph{首项、高阶项与 $\pi\sim\Li$}
\begin{enumerate}
  \item $x/\log x$ 是 $\Li$ 渐近展开的首项，见~\eqref{eq:Li-parts-1}、\eqref{eq:Li-asymptotic}。
  \item 素数定理断言 $\pi(x)\sim\Li(x)$，从而也有 $\pi\sim x/\log x$。
    在证明之前，表~\ref{tab:pix} 与图~\ref{fig:pix_compare} 给出数值对照；
    证明提要见第~\ref{ch:prime_number_theorem}~章，细节见附录~\ref{app:pnt_estimates}。
  \item Legendre 分母修正与 $\Li$ 高阶项的关系。
    由~\eqref{eq:Li-asymptotic}，
    \[
      \Li(x)
      \sim
      \frac{x}{\log x}
      +\frac{x}{(\log x)^{2}}
      +\frac{2x}{(\log x)^{3}}
      +\frac{6x}{(\log x)^{4}}
      +\cdots.
    \]
    首项即 $x/\log x$；其后各项统称高阶项。
    Legendre 的经验近似
    \[
      \pi(x)\;\approx\;\frac{x}{\log x-A}
      \qquad (A\approx 1.08366)
    \]
    将分母写成 $\log x-A$，展开得
    \begin{equation}
      \frac{x}{\log x-A}
      =\frac{x}{\log x}\Bigl(1-\frac{A}{\log x}\Bigr)^{-1}
      =\frac{x}{\log x}
      +\frac{A\,x}{(\log x)^{2}}
      +\frac{A^{2}\,x}{(\log x)^{3}}
      +\cdots,
      \label{eq:legendre-denom-expand}
    \end{equation}
    即以单一拟合常数 $A$ 同时影响首项之后各项，相当于用一条经验近似去逼近~\eqref{eq:Li-asymptotic} 中的多阶贡献；
    $A$ 由数值拟合得到，与~\eqref{eq:Li-asymptotic} 中的系数 $1,2,6,\ldots$ 一般并不相同。
    最简情形 $A=0$ 即 $x/\log x$（图~\ref{fig:pix_compare} 中的绿线）。
    Gauss 的 $\Li$ 则以积分一次写出~\eqref{eq:Li-asymptotic} 的光滑和；
    后文显式公式中，对数积分才作为与零点级数相对的主项出现。
\end{enumerate}

\begin{table}[htbp]
\centering
\caption{$\pi(x)$ 与 $x/\log x$、$\Li(x)$ 的数值对照。}
\label{tab:pix}
\renewcommand{\arraystretch}{1.3}
\begin{tabular}{r r r r r}
\toprule
$x$ & $\pi(x)$ & $\left\lfloor\dfrac{x}{\log x}\right\rfloor$ & $\left\lfloor\Li(x)\right\rfloor$ & $\pi(x) - \left\lfloor\Li(x)\right\rfloor$ \\[6pt]
\midrule
$10$ & $4$ & $4$ & $5$ & $-1$ \\
$10^2$ & $25$ & $21$ & $29$ & $-4$ \\
$10^3$ & $168$ & $144$ & $176$ & $-8$ \\
$10^4$ & $1\,229$ & $1\,085$ & $1\,245$ & $-16$ \\
$10^5$ & $9\,592$ & $8\,685$ & $9\,628$ & $-36$ \\
$10^6$ & $78\,498$ & $72\,382$ & $78\,626$ & $-128$ \\
$10^7$ & $664\,579$ & $620\,420$ & $664\,917$ & $-338$ \\
$10^8$ & $5\,761\,455$ & $5\,428\,681$ & $5\,762\,208$ & $-753$ \\
$10^9$ & $50\,847\,534$ & $48\,254\,942$ & $50\,849\,233$ & $-1\,699$ \\
$10^{10}$ & $455\,052\,511$ & $434\,294\,481$ & $455\,055\,613$ & $-3\,102$ \\
\bottomrule
\end{tabular}
\end{table}

表~\ref{tab:pix} 显示：$\Li$ 的相对误差随 $x$ 增大而迅速减小，且在表中取值范围内 $\pi(x)-\Li(x)<0$。
Littlewood（1914）证明 $\pi-\Li$（及与 $\li$ 的差）变号无穷多次；
使 $\pi(x)>\Li(x)$ 首次成立的 $x$ 极大
（上界与 Skewes 数见第~\ref{ch:elementary_number_theoretic_functions}~章第~\ref{subsec:Li_pi_asymp}~节）。

\begin{figure}[htbp]
\centering
\begin{tikzpicture}
\begin{axis}[
  width=0.95\textwidth,
  height=8.2cm,
  xlabel={$x$},
  ylabel={},
  xmin=2, xmax=100,
  ymin=0, ymax=32,
  xtick={10,20,...,100},
  ytick={0,5,...,30},
  grid=major,
  grid style={line width=0.3pt, draw=gray!30},
  axis lines=left,
  legend pos=north west,
  legend style={font=\small, cells={anchor=west}, row sep=2pt},
  legend cell align=left,
]
% 中等区间：三条曲线可分；大 $x$ 时 $|\pi-\Li|/\pi$ 很小，宏观图上会几乎重合
\addplot[blue!75!black, very thick, const plot, no marks] coordinates {
  (1,0) (2,1) (3,2) (5,3) (7,4) (11,5) (13,6) (17,7) (19,8) (23,9)
  (29,10) (31,11) (37,12) (41,13) (43,14) (47,15) (53,16) (59,17) (61,18)
  (67,19) (71,20) (73,21) (79,22) (83,23) (89,24) (97,25) (100,25)
};
\addlegendentry{$\pi(x)$}
\addplot[red!80!black, thick, no marks] coordinates {
  (2,0.00)(3,1.12)(4,1.92)(5,2.59)(6,3.18)(7,3.71)(8,4.21)(9,4.68)(10,5.12)
  (11,5.55)(12,5.96)(13,6.35)(14,6.74)(15,7.11)(16,7.47)(17,7.83)(18,8.18)
  (19,8.52)(20,8.86)(21,9.19)(22,9.52)(23,9.84)(24,10.16)(25,10.47)(26,10.78)
  (27,11.08)(28,11.38)(29,11.68)(30,11.98)(32,12.56)(34,13.13)(36,13.69)
  (38,14.25)(40,14.79)(42,15.33)(44,15.86)(46,16.39)(48,16.91)(50,17.42)
  (55,18.69)(60,19.92)(65,21.13)(70,22.32)(75,23.48)(80,24.63)(85,25.77)
  (90,26.88)(95,27.99)(100,29.08)
};
\addlegendentry{$\Li(x)$（Gauss 逼近）}
\addplot[green!55!black, very thick, solid, domain=2.8:100, samples=80] {x/ln(x)};
\addlegendentry{$x/\log x$（Legendre 近似）}
\end{axis}
\end{tikzpicture}
\caption{区间 $2\le x\le 100$ 上 $\pi(x)$、$\Li(x)$ 与 $x/\log x$ 的比较。
$\Li$ 介于 $\pi$ 与 $x/\log x$ 之间且更接近 $\pi$，$x/\log x$ 整体偏低
（见~\eqref{eq:Li-parts-1}）。$x$ 较大时相对误差变小，曲线在粗尺度上几乎重合，
故取中等区间以显示绝对差。}
\label{fig:pix_compare}
\end{figure}

%===========================================================
\section{素数定理的表述}\label{subsec:pnt_intro}
%===========================================================

上文的 $\Li$ 与 $x/\log x$ 原为经验公式；素数定理断言它们与 $\pi(x)$ 渐近等价。
本节仅给出定理表述，证明见后文。

\begin{theorem}{素数定理，1896}{}
\label{thm:pnt}
当 $x \to \infty$ 时，
\[
\pi(x) \sim \Li(x) \sim \frac{x}{\log x}.
\]
即 $\lim_{x\to\infty} \frac{\pi(x)}{\Li(x)} = 1$。
\end{theorem}

Hadamard 与 de~la~Vallée-Poussin 于 1896 年独立证明；关键步骤是证明 $\zeta(s)$ 在 $\operatorname{Re}s=1$（$s\neq 1$）上无零点。
证明提要见第~\ref{ch:prime_number_theorem}~章，细节见附录~\ref{app:pnt_estimates}。
$\Li(x)>x/\log x$（当 $x$ 充分大）已由上文余项说明；
$\pi-\Li$ 的形状与量级见下一节及第~\ref{ch:riemann_explicit_formula}~章。
（记号：素数定理与数值比较用实 $\Li$；显式公式用 $\li(x^{\rho})$；
分工见第~\ref{ch:logint_Ei}~章。）

由此可见：在解析方法出现之前，数学家虽已掌握 Legendre 型与 Gauss 对数积分等经验公式，
对如何证明素数定理却长期缺乏有效路径。
Riemann 将素数分布引入复分析框架，建立了新的研究范式，
为上述 1896 年证明提供了理论基础，问题始得解决。
黎曼猜想至今未决，数学界仍在探索足以推进它的新思想、新工具或新范式。

%===========================================================
\section{从经验公式到严格理论}\label{sec:riemann_hypothesis_intro}
%===========================================================

古典公式原由数值观察提出；经解析方法方可成为可严格证明的数学定理。
这一跃迁的线索可概括为
\[
\underbrace{\frac{x}{\log x}}_{\text{Legendre 型}}
\;\longrightarrow\;
\underbrace{\Li(x)}_{\text{Gauss 对数积分}}
\;\longrightarrow\;
\underbrace{\text{黎曼显式公式}}_{\text{$\li/\Li$ 进入公式}}
\;\longrightarrow\;
\underbrace{\pi\sim\Li}_{\text{PNT，1896}}.
\]
关键一步在于：黎曼利用复分析方法，推导出素数分布的显式公式；
原先作为经验逼近的对数积分 $\Li$（及其同族的 $\li$）出现在显式公式中，
从而与后来可严格证明的素数定理理论建立了数学联系——
由经验估计得到的公式不再停留在数值拟合上，而被写入一条可由复分析严格证明的数学表达式之中。
素数定理 $\pi\sim\Li$ 表明，Gauss 的直觉是敏锐的；
对数积分 $\Li$ 作为逼近主项出现在黎曼显式公式中，表明了 $\Li$ 与素数分布的底层逻辑关系。
图~\ref{fig:zeta_zeros_primes} 是黎曼推导显式公式的主要节点与路径的简单示意；
以下四小节对图中四个节点依次作简单介绍。
完整证明见第~\ref{ch:prime_number_theorem}、\ref{ch:riemann_explicit_formula}~章。

\begin{figure}[htbp]
\centering
\resizebox{0.92\textwidth}{!}{%
\begin{tikzpicture}[
  node distance=1.8cm,
  box/.style={draw, rounded corners=4pt, minimum width=2.8cm,
              minimum height=0.95cm, align=center, font=\small,
              fill=blue!6, draw=blue!40},
  arrow/.style={-{Stealth[length=6pt]}, thick, blue!60}
]
  \node[box] (A) {Euler 乘积};
  \node[box, right=of A] (B) {非平凡零点};
  \node[box, right=of B] (C) {显式公式};
  \node[box, right=of C] (D) {$\pi\sim\Li$};
  \draw[arrow] (A) -- (B);
  \draw[arrow] (B) -- (C);
  \draw[arrow] (C) -- (D);
\end{tikzpicture}%
}
\caption{黎曼推导显式公式的主要节点与路径（示意）。}
\label{fig:zeta_zeros_primes}
\end{figure}

\subsection{Euler 乘积}
\label{subsec:euler_product_complex}

对 $\operatorname{Re}(s)>1$，Euler 给出
\begin{equation}
  \sum_{n=1}^{\infty}\frac{1}{n^{s}}
  =\prod_{p}\frac{1}{1-p^{-s}},
  \label{thm:euler-product-ch1}
\end{equation}
乘积取遍素数。它把「全体素数」写进一个依赖复变元 $s$ 的公式，是解析进路的入口。

\subsection{非平凡零点与黎曼猜想}
\label{subsec:zeros_rh_intro}

上述公式经延拓后，在复平面上除极点外的零点中，落在带状区域 $0<\operatorname{Re}s<1$ 内的称为\textbf{非平凡零点}。
它们将进入计数公式的修正项。

\subsection{显式公式（形状）}
\label{subsec:explicit_formula_intro}

黎曼显式公式把素数计数写成：\textbf{以对数积分为主的光滑部分}，加上\textbf{由非平凡零点给出的修正}。
经验上的 $\li/\Li$ 由此进入可论证的恒等式，并与 $\pi\sim\Li$ 衔接。

\subsection{素数定理 \texorpdfstring{$\pi\sim\Li$}{pi ~ Li}}
\label{subsec:pnt_as_endpoint}

图~\ref{fig:zeta_zeros_primes} 的末节点是素数定理
\[
  \pi(x)\sim\Li(x)
  \qquad (x\to\infty),
\]
亦即 $\pi(x)/\Li(x)\to 1$。
经验上的 $\Li$ 经显式公式进入严格理论之后，渐近等价 $\pi\sim\Li$ 成为可证明的定理（1896）。

\begin{center}
\fbox{\parbox{0.92\textwidth}{
\centering
\vspace{0.2em}
\textbf{第一章要点}\\[6pt]
\renewcommand{\arraystretch}{1.5}
\begin{tabular}{lp{9.2cm}}
\hline
计数对象 & $\pi(x)=\#\{p\le x\}=\sum_{p\le x}1$ \\
经验公式 & $x/\log x$；$\Li(x)=\int_2^x\mathrm{d}t/\log t$ \\
$\Li$ 展开 & $x/\log x$ 为首项；$\Li\sim\sum_{k\ge1}(k-1)!\,x/(\log x)^{k}$ \\
Legendre 型近似 & $\pi(x)\approx x/(\log x-A)$，$A\approx 1.08366$；$A=0$ 即 $x/\log x$ \\
素数定理 & $\pi\sim\Li\sim x/\log x$（表述；证明在后文） \\
解析路线（主要线索） & Euler 乘积 $\to$ 零点 $\to$ 显式公式（$\li/\Li$ 入主项）$\to$ PNT \\
\hline
\end{tabular}
\vspace{0.2em}
}}
\end{center}

%===========================================================
\section*{本章小结与后续展望}
\addcontentsline{toc}{section}{本章小结与后续展望}
%===========================================================

本章引入素数分布问题：$\pi(x)$、古典经验公式、素数定理的表述，
以及通向后文的解析路线主要线索（Euler 乘积—零点—显式公式）。
记号见《符号与约定》；$\pi,J,\psi,\mu,M,\Lambda,\Li$ 的系统定义见第~\ref{ch:elementary_number_theoretic_functions}~章。

全书各章的前提关系见图~\ref{fig:roadmap}；箭头由前提指向后继。
读图约定与主要前提边列于图后（与图同页衔接，避免长题注被挤到他页）。

\begin{figure}[!ht]
\centering
\resizebox{!}{0.72\textheight}{%
\begin{tikzpicture}[
  font=\small,
  >=Stealth,
  ch/.style={
    draw=MainBlue!55!black,
    fill=LightBlue!30,
    rounded corners=3pt,
    align=center,
    inner sep=6pt,
    text width=5.6cm,
    minimum height=1.05cm,
    font=\small
  },
  hub/.style={
    draw=DarkGreen!65!black,
    fill=LightGreen!40,
    rounded corners=3pt,
    align=center,
    inner sep=6pt,
    text width=5.6cm,
    minimum height=1.05cm,
    font=\small
  },
  key/.style={
    draw=Gold!75!black,
    fill=LightGold!55,
    rounded corners=3pt,
    align=center,
    inner sep=6pt,
    text width=5.6cm,
    minimum height=1.05cm,
    font=\small
  },
  arr/.style={->, thick, MainBlue!75!black},
  wire/.style={thick, MainBlue!75!black},
  % 电路图约定：● = 连线在此汇合（有关联）；交叉无点 = 仅几何掠过、无关联
  jdot/.style={circle, fill=MainBlue!90!black, draw=none,
    inner sep=0pt, minimum size=3.4pt},
]
  % ========== 全部章框纵向一列 ==========
  % 第1章略上移，加大与第2章间距，保证 1→2 箭头杆完整可见
  \node[ch, fill=LightGold!45, draw=Gold!70!black] (ch1) at (0,0.55)
    {第 1 章（本章）：经验公式 · PNT 表述 · 解析理论线索};
  \node[ch] (ch2) at (0,-1.55)
    {第 2 章：数论函数 $\pi,J,\psi,\mu,M,\Lambda,\Li$};
  \node[ch] (ch3) at (0,-3.1)
    {第 3 章：复变函数简介};
  \node[ch] (ch4) at (0,-4.65)
    {第 4 章：积分变换（Mellin / Perron）};
  \node[ch] (ch5) at (0,-6.2)
    {第 5 章：$\zeta$ 初探（$\operatorname{Re}s>1$）};
  \node[ch] (ch6) at (0,-7.75)
    {第 6 章：离散部分和与 Perron 积分};
  \node[ch] (ch7) at (0,-9.3)
    {第 7 章：$\Gamma$ 函数};
  \node[ch] (ch8) at (0,-10.85)
    {第 8 章：解析延拓的定义与唯一性};
  \node[hub] (ch9) at (0,-12.4)
    {第 9 章：$\zeta$ 延拓与函数方程};
  \node[ch] (ch10) at (0,-13.95)
    {第 10 章：$\xi$ 与 Hadamard};
  \node[ch] (ch11) at (0,-15.5)
    {第 11 章：零点几何与 $N(T)$};
  \node[hub] (ch12) at (0,-17.05)
    {第 12 章：素数定理（证明提要）};
  \node[hub] (ch13) at (0,-18.6)
    {第 13 章：$\Li$、$\li$ 与 $\Ei$};
  \node[key] (ch14) at (0,-20.15)
    {第 14 章：黎曼显式公式};
  \node[ch, fill=white] (ch15) at (0,-21.7)
    {第 15--20 章：谱系图示 · 逼近 · 数值 · 几何示意 · 原稿 · 前沿};

  % ========== 主链（框间直连）==========
  \draw[arr] (ch1.south) -- (ch2.north);  % 1→2 画全（锚点对接）
  \draw[arr] (ch3) -- (ch4);
  \draw[arr] (ch5) -- (ch6);
  \draw[arr] (ch9) -- (ch10);
  \draw[arr] (ch10) -- (ch11);
  \draw[arr] (ch11) -- (ch12);
  \draw[arr] (ch12) -- (ch13);
  \draw[arr] (ch13) -- (ch14);
  \draw[arr] (ch14) -- (ch15);

  % ========== 左侧总线：第2章 → 第5、6、14（有关联的 T 形汇合加点）==========
  % 竖廊 x；仅在「接入总线 / 从总线接出」处画 ●
  \coordinate (Lx) at (-4.15,0);
  \coordinate (L2) at (ch2.west-|Lx);
  \coordinate (L5) at (ch5.west-|Lx);
  \coordinate (L6) at (ch6.west-|Lx);
  \coordinate (L14) at (ch14.west-|Lx);
  \draw[wire] (ch2.west) -- (L2);
  \draw[wire] (L2) -- (L5) -- (L6) -- (L14);
  \draw[arr] (L5) -- (ch5.west);
  \draw[arr] (L6) -- (ch6.west);
  \draw[arr] (L14) -- (ch14.west);
  % L2、L14 为单线折角，不设点；L5/L6 为总线分出的 T 形汇合
  \node[jdot] at (L5) {};
  \node[jdot] at (L6) {};

  % 第5→第9：独立竖廊（与第2章总线分离，交叉处无点 = 无关联）
  \coordinate (L59x) at (-5.05,0);
  \coordinate (L59a) at (ch5.west-|L59x);
  \coordinate (L59b) at (ch9.west-|L59x);
  \draw[wire] (ch5.west) -- (L59a);
  \draw[wire] (L59a) -- (L59b);
  \draw[arr] (L59b) -- (ch9.west);
  % 与左侧总线若几何接近，不设汇合点（无 5–总线 以外的关联）

  % 第10→第12：单条折线，直角转角不是交叉汇合，不设点
  \coordinate (L1012x) at (-3.35,0);
  \draw[arr] (ch10.west) -- (ch10.west-|L1012x) -- (ch12.west-|L1012x) -- (ch12.west);

  % ========== 右侧 ==========
  % 约定：● 仅多线真正汇合（T 形）；单线折角、框边出线中点不设点
  \coordinate (R9x)  at (4.50,0);   % →9 总线
  \coordinate (R14x) at (5.45,0);   % →14 总线

  % ----- 3,4,7,8 → 9 -----
  \coordinate (R3)  at (ch3.east-|R9x);
  \coordinate (R4a) at (ch4.east-|R9x);
  \coordinate (R7)  at (ch7.east-|R9x);
  \coordinate (R8)  at (ch8.east-|R9x);
  \coordinate (R9)  at (ch9.east-|R9x);
  \draw[wire] (ch3.east) -- (R3);
  \draw[wire] (ch4.east) -- (R4a);
  \draw[wire] (ch7.east) -- (R7);
  \draw[wire] (ch8.east) -- (R8);
  \draw[wire] (R3) -- (R4a) -- (R7) -- (R8) -- (R9);
  \draw[arr] (R9) -- (ch9.east);
  % 仅竖廊中段的 T 形汇合（有上、下竖线 + 水平支线）；端点折角 R3/R9 不设点
  \foreach \P in {R4a,R7,R8} \node[jdot] at (\P) {};

  % ----- 4→14、6→14：共竖廊（不再在框旁另画 4→6，以免与竖廊/主链重复）-----
  % R4b 单线折角不设点；R6b 为 6 与竖廊的 T 汇合；R14 接到第 14 章
  \coordinate (R4b) at (ch4.east-|R14x);
  \coordinate (R6b) at (ch6.east-|R14x);
  \coordinate (R14) at (ch14.east-|R14x);
  \draw[wire] (ch4.east) -- (R4b) -- (R6b) -- (R14);
  \draw[wire] (ch6.east) -- (R6b);
  \draw[arr] (R14) -- (ch14.east);
  % R4b、R14 为直角折角，不设点；仅 R6b 为 T 形汇合
  \node[jdot] at (R6b) {};

\end{tikzpicture}%
}
\caption{全书各章的前提关系；箭头由前提指向后继。}
\label{fig:roadmap}
\end{figure}

\noindent\textbf{读图约定}
（同电路图）：\textbf{实心小圆点}标在连线\textbf{真正汇合}处，表示有前提关联；
折线交叉若\textbf{无圆点}，仅为走线掠过，\textbf{不}表示两章在交点处有关联。
例如第~2~章经左侧总线连至第~5、6、14~章（汇合处有点）；
第~3、4、7、8~章经右侧总线汇入第~9~章；第~4、6~章经另一竖廊汇入第~14~章——
第~3~章与第~14~章无直接前提边，其间若有线条交叉亦无圆点。

\noindent\textbf{主要前提边}
主链：第~12$\to$13$\to$14；
另有 $2\to5$、$2\to6$、$2\to14$；$3\to4$、$3\to9$；
$4\to9$、$4\to14$；$5\to6$、$5\to9$；$6\to14$；
$7\to9$、$8\to9$；$10\to12$ 等
（第~4$\to$6 由第~5$\to$6 主链及第~4、6 共入第~14~章之关系体现，图中不在框旁另画）。
